\documentclass[UTF8,a4paper,12pt]{ctexart}

\usepackage[margin=2.2cm]{geometry}
\usepackage{hyperref}
\usepackage{booktabs}
\usepackage{longtable}
\usepackage{array}
\usepackage{enumitem}
\usepackage{xcolor}
\usepackage{amssymb}

\hypersetup{
  colorlinks=true,
  linkcolor=blue,
  urlcolor=blue
}

\setlist[itemize]{itemsep=2pt,topsep=2pt,parsep=0pt,leftmargin=2em}
\setlist[enumerate]{itemsep=2pt,topsep=2pt,parsep=0pt,leftmargin=2.2em}

\title{北京市``两节两营''科技冬令营课程方案（口令攻防模块·两天）\\
\large 奇安信：新一代网络安全领军者\\
\normalsize \textit{To build is to understand; To break is to secure.}}
\author{口令攻防模块负责人（草案）}
\date{\today}

\begin{document}
\maketitle

\section*{模块定位与原则}
\begin{itemize}
  \item \textbf{面向对象：}初中为主（可混编高中生担任小组助教/记录员）；学生默认\textbf{无密码学基础}、无攻防经验。
  \item \textbf{模块位置：}7天课程中的第3--4天（两天），承接第1--2天的``个人建站''成果（每组已有可访问的练习站点/登录页）。
  \item \textbf{核心主线：}从\textbf{在线暴力尝试}的直观风险出发，理解验证码/人机验证/限速/锁定的意义；第2天扩展到\textbf{个性化口令猜测}（TarGuess-i，PCFG-based），最后\textbf{简提}LLM带来的新风险。
  \item \textbf{安全边界：}所有攻击演示仅限\textbf{本地/课堂内网靶场}，使用\textbf{虚构账号与口令空间}（例如4位PIN或极小字母表），\textbf{不}教授/不布置通用化的外网攻击。
  \item \textbf{评价方式：}以\textbf{过程性考核}为主（理解、记录、解释、复现与改造），结果性考核为辅（功能是否生效、对比是否完整）。
\end{itemize}

\section{学习目标（Learning Outcomes）}
\subsection*{知识目标}
\begin{itemize}
  \item 说清楚HTTPS/证书的作用：\textbf{身份认证}与\textbf{传输加密}分别解决什么问题。
  \item 理解最小密码学概念：\textbf{哈希}（不可逆指纹）、\textbf{加盐}（避免相同口令产生相同指纹）、\textbf{慢哈希}（提升离线猜测成本）。
  \item 理解在线口令猜测的核心特征：\textbf{短时间大量失败请求}与\textbf{搜索空间}。
  \item 理解PCFG/TarGuess-i的直观思想：\textbf{结构/模式+概率排序}；知道它与``纯暴力''的差别。
\end{itemize}

\subsection*{技能目标}
\begin{itemize}
  \item 能运行给定的演示脚本/工具并解读输出（尝试次数、速率、命中位置）。
  \item 能对练习网站实现至少1项防护改造（验证码/人机验证概念、限速、失败锁定、延迟退避、简化二步验证等）并完成复测对比。
  \item 能基于人物画像列出Top20口令候选并总结常见模式，再与TarGuess-i输出进行对比分析。
\end{itemize}

\subsection*{态度与意识}
\begin{itemize}
  \item 形成``\textbf{理解攻击才能设计防御}''的工程思维；能在安全与可用性之间做解释性取舍。
  \item 形成对``公开个人信息$\rightarrow$个性化猜测风险''的警觉，能提出可执行的个人与系统防护建议。
\end{itemize}

\section{课程组织与资源}
\subsection{组织方式（建议）}
\begin{itemize}
  \item 班级300人建议按\textbf{4--6人/组}分组；每组设置：组长、开发、测试、记录员（高中生可优先担任测试/记录/助教）。
  \item 统一发放``练习站点项目''（Flask示例）、人物卡、实验记录表（纸质或在线表单）。
\end{itemize}

\subsection{软硬件资源}
\begin{itemize}
  \item 学生端：浏览器、Python环境（建议3.10+）、IDE（任选：VS Code / PyCharm Edu / Thonny），命令行基础操作（进入目录、运行脚本）。
  \item 教师端：投影、可控靶场网络（本地/内网）、预先准备的演示脚本（在线尝试）、TarGuess-i运行环境与示例数据。
  \item 练习站点：提供带登录页的Flask小站（例如本项目），可快速启动与修改。
\end{itemize}

\section{两天总流程与时间建议}
\subsection*{总览（可按实际压缩/拉伸）}
\begin{longtable}{@{}p{0.12\textwidth}p{0.18\textwidth}p{0.64\textwidth}@{}}
\toprule
\textbf{日程} & \textbf{环节} & \textbf{内容与产出} \\
\midrule
\endhead
Day 1 & 环节0：基础科普 & HTTPS/证书、最小密码学概念、口令策略；产出：口令规则卡（小组版）\\
Day 1 & 环节1：在线爆破与防护 & 教师演示在线暴力尝试；学生完成验证码/人机验证概念、限速/锁定等改造并复测；产出：前后对比记录 \\
Day 2 & 环节2：个性化猜测 & 人物画像猜口令Top20；总结模式；运行TarGuess-i并对比命中排名；产出：模式清单+对比分析 \\
Day 2 & 扩展：LLM与离线 & 10--20分钟：LLM猜测趋势、离线攻击概念（闲谈）；产出：2条个人建议+2条系统建议 \\
\bottomrule
\end{longtable}

\section{Day 1：基础知识 + 在线爆破与防护改造}

\subsection{环节0（科普）：HTTPS/证书与最小密码学}
\subsubsection*{基础知识要点（讲义结构建议）}
\begin{itemize}
  \item \textbf{问题1：口令在路上会发生什么？}\\
  浏览器$\rightarrow$网络$\rightarrow$服务器。没有HTTPS时可能被\textbf{窃听}（看见你输入）或\textbf{篡改}（把你导向假网站）。
  \item \textbf{证书是什么？}\\
  证书像``网站身份证''：让浏览器确认``你连的是对的那个网站''（身份认证）。
  \item \textbf{TLS/HTTPS解决什么？}\\
  TLS建立加密通道：解决\textbf{传输机密性}与\textbf{完整性}，让旁路的人看不懂也改不了。
  \item \textbf{问题2：服务器怎么保存口令才不危险？}\\
  口令不能明文保存。理想做法：保存\textbf{哈希}（不可逆）、加\textbf{盐}（每人不同）、使用\textbf{慢哈希}（让离线猜测变慢）。
  \item \textbf{问题3：好口令长什么样？}\\
  长度优先（建议$\ge 12$字符）、不绑定明显个人信息、不复用、重要账户开启MFA。
\end{itemize}

\subsubsection*{小活动0A：``看一眼证书''（5--10分钟，非必需）}
\paragraph{原理}
浏览器会展示HTTPS连接的证书信息：签发给谁、由谁签发、有效期等。证书异常会出现警告。
\paragraph{步骤}
\begin{enumerate}
  \item 打开任意HTTPS站点（可用校内门户或示例站），点击地址栏锁形图标查看证书。
  \item 观察证书的\textbf{域名}、\textbf{有效期}、\textbf{签发机构}。
  \item 讨论：如果证书域名不匹配或过期，可能意味着什么风险？
\end{enumerate}
\paragraph{过程考核点}
学生能用一句话解释：证书主要解决``\underline{\hspace{6em}}''问题（身份认证）。

\subsubsection*{小活动0B：口令强度与``搜索空间''（10--15分钟）}
\paragraph{原理}
口令越长、可用字符越多，暴力尝试需要的组合数就越大。课堂只需用直觉理解数量级变化（不要求计算对数）。
\paragraph{步骤}
\begin{enumerate}
  \item 教师给出三种口令规则（示例）：4位数字PIN、8位小写字母、12位混合字符。
  \item 学生只做``更大/更小''比较：哪种最容易被试出来？为什么？
  \item 小结：\textbf{长度优先}通常比``加几个符号''更关键。
\end{enumerate}
\paragraph{过程考核点}
学生能解释：把口令长度从8提高到12，在线爆破难度会\underline{\hspace{6em}}（显著上升）。

\subsection{环节1（主实验）：在线口令爆破演示 + 防护改造}

\subsubsection{实验1：在线暴力尝试（教师演示为主，学生观察记录）}
\paragraph{实验目的}
让学生理解``在线爆破''的本质：对登录接口做大量重复尝试；理解\textbf{搜索空间}与\textbf{尝试速率}决定风险。

\paragraph{实验原理（讲给学生的版本）}
把口令看作一把钥匙。如果钥匙只有4位数字（0000--9999），攻击者可以按顺序试：试的速度越快、口令越短，就越容易被撞中。\\
在线爆破的特征是：同一账号（或同一IP）出现\textbf{短时间大量失败登录}。

\paragraph{实验准备（教师）}
\begin{itemize}
  \item 启动练习站点（本地/内网），准备一个演示账号（虚构），口令设置为\textbf{小空间}（例如4位PIN或极简字符集）。
  \item 准备演示脚本：向\texttt{/login}连续发送请求（受控速率、受控总次数）。
  \item 明确课堂规则：禁止对外网、禁止对同学账号；只对老师提供的靶站与靶账号测试。
\end{itemize}

\paragraph{演示脚本安全约束（建议写进课堂规则）}
\begin{itemize}
  \item 仅允许访问课堂靶站的域名/IP白名单；禁止任意URL输入。
  \item 速率上限：例如$\le 5$次/秒（保证网络与服务稳定）。
  \item 总次数上限：例如$\le 2000$次；达到上限必须自动停止。
  \item 口令空间上限：仅允许PIN或极小字符集，避免把课堂变成``真实破解训练''。
\end{itemize}

\paragraph{步骤（课堂执行）}
\begin{enumerate}
  \item 教师投影展示：正常手动登录一次（说明登录入口与参数）。
  \item 教师运行演示脚本：以小口令空间进行尝试，直到命中或达到次数上限。
  \item 学生填写观察表：记录尝试总数、用时、平均速率、命中位置（第几次）。
  \item 全班讨论：把口令长度从4位提升到6位，风险变化会是什么量级？
\end{enumerate}

\paragraph{学生记录表（实验1）}
\begin{longtable}{@{}p{0.22\textwidth}p{0.72\textwidth}@{}}
\toprule
\textbf{字段} & \textbf{填写内容} \\
\midrule
\endhead
口令空间假设 & 例如：0000--9999（共10000种） \\
尝试上限 & 例如：500次 / 2000次 \\
实际用时 & \underline{\hspace{12em}} \\
平均速率（估算） & \underline{\hspace{12em}} 次/秒 \\
是否命中 & 是 / 否 \\
命中位置（第几次） & \underline{\hspace{12em}} \\
风险一句话结论 & \underline{\hspace{20em}} \\
\bottomrule
\end{longtable}

\subsubsection{实验2：防护改造（小组动手）}
\paragraph{实验目的}
把``能被爆破''的登录系统改造成``不容易被自动化尝试''：让攻击\textbf{变慢}、\textbf{变贵}、\textbf{变不可持续}。

\paragraph{实现提示（以Flask练习站为例，面向学生的任务拆解）}
以下是推荐的实现拆解（不要求一次做全，可按小组能力选做）：
\begin{itemize}
  \item \textbf{状态记录：}为``账号''或``IP''维护登录失败计数与最后失败时间（可先用内存字典，强调课堂演示即可）。
  \item \textbf{阈值策略：}定义阈值$N$与锁定时长$T$；超过阈值则拒绝登录并提示剩余时间。
  \item \textbf{挑战机制：}失败次数达到$k$后才出现验证码/算术题（降低正常用户负担）。
  \item \textbf{可观测性：}把登录事件写入简单日志（时间、账号、成功/失败、触发了什么限制），用于复盘。
\end{itemize}
若使用本仓库示例项目，可直接在\texttt{app.py}与\texttt{templates/login.html}中完成最小改造（教师可提供起始模板文件，避免学生从零写）。

\paragraph{实验原理（分层解释）}
\begin{itemize}
  \item \textbf{验证码/人机验证：}让每次尝试需要人类参与或更高成本的识别，从而抑制自动化。
  \item \textbf{限速/延迟：}每次失败增加等待，让单位时间尝试次数下降。
  \item \textbf{失败锁定：}失败次数超过阈值后短暂锁定账号或要求更强验证。
  \item \textbf{（可选）简化二步验证：}即使口令被猜中，还需要第二个因素才能登录成功。
\end{itemize}

\paragraph{任务分级（每组选做，推荐至少完成A）}
\begin{enumerate}[label=\textbf{\Alph*.}]
  \item \textbf{失败次数阈值 + 临时锁定（推荐）}\\
  例如：同一账号连续失败3次，锁定60秒，期间不再处理登录尝试。
  \item \textbf{失败后延迟/退避（推荐）}\\
  例如：每次失败增加1秒等待，最多增加到5秒（演示``爆破变慢''）。
  \item \textbf{验证码（推荐）}\\
  使用课堂可实现的简单形式：算术题/拼图验证码/一次性图片码（不追求工业强度，追求原理）。
  \item \textbf{人机验证/Cloudflare概念展示（讲解型）}\\
  说明网关层如何做风控：指纹、速率限制、挑战页面。若课堂环境无法真正接入Cloudflare，可用``模拟网关''规则展示思想。
  \item \textbf{（挑战）简化二步验证}\\
  登录口令正确后再输入一次性码（教师发放/小组共享密钥生成），或回答预设安全问题（强调缺点：安全问题易被猜）。
\end{enumerate}

\paragraph{步骤（学生）}
\begin{enumerate}
  \item 阅读项目的登录流程：前端表单$\rightarrow$后端\texttt{/login}校验$\rightarrow$写入session。
  \item 选择任务A--E之一，完成代码改造与自测。
  \item 填写``防护设计说明''：\textbf{我在对抗哪种特征？}（高频失败/自动化/重复尝试）
  \item 由教师或助教用同一演示脚本进行复测，学生记录前后对比。
\end{enumerate}

\paragraph{小组提交物（实验2）}
\begin{itemize}
  \item 防护点说明（1页以内）：做了什么、为什么有效、会带来什么副作用（可用性）。
  \item 复测对比表：改造前/后尝试速率、是否被锁定/是否需要验证码、成功率变化。
\end{itemize}

\paragraph{复测对比表（实验2）}
\begin{longtable}{@{}p{0.18\textwidth}p{0.36\textwidth}p{0.36\textwidth}@{}}
\toprule
\textbf{指标} & \textbf{改造前} & \textbf{改造后} \\
\midrule
\endhead
每秒尝试次数（估算） & \underline{\hspace{10em}} & \underline{\hspace{10em}} \\
达到阈值是否触发限制 & 否 / 是（\underline{\hspace{6em}}） & 否 / 是（\underline{\hspace{6em}}） \\
是否需要验证码/挑战 & 否 / 是（类型：\underline{\hspace{6em}}） & 否 / 是（类型：\underline{\hspace{6em}}） \\
攻击是否还能自动化 & 是 / 否（理由：\underline{\hspace{6em}}） & 是 / 否（理由：\underline{\hspace{6em}}） \\
可用性影响 & \underline{\hspace{18em}} & \underline{\hspace{18em}} \\
\bottomrule
\end{longtable}

\subsubsection*{Day 1 过程考核表（教师/助教）}
\begin{longtable}{@{}p{0.28\textwidth}p{0.10\textwidth}p{0.54\textwidth}@{}}
\toprule
\textbf{考核点} & \textbf{达成} & \textbf{备注} \\
\midrule
\endhead
能解释证书/TLS解决的两个问题（身份、加密） & $\square$ & \\
能用自己的话解释``哈希+盐''的意义（概念级） & $\square$ & \\
完成实验1观察记录（字段完整） & $\square$ & \\
完成至少一项防护改造（A--E）并可演示 & $\square$ & \\
能说明所做防护针对的攻击特征与副作用 & $\square$ & \\
\bottomrule
\end{longtable}

\subsection*{Day 1 小测（可选，5题，10分钟）}
\begin{enumerate}
  \item 证书主要用来解决什么问题？（A. 让网页更好看 B. 身份认证 C. 提升网速）
  \item HTTPS主要保护的是：A. 传输中的机密性/完整性 B. 电脑不被病毒感染 C. 账号永不泄露
  \item 为什么不能把口令明文存数据库？请用一句话说明风险。
  \item 在线爆破的两个关键量是什么？（提示：空间与速度）
  \item 写出一种你能想到的登录防护策略，并写一句话解释它如何降低爆破成功率。
\end{enumerate}

\section{Day 2：人物画像猜口令 + TarGuess-i（PCFG）+ 简提LLM}

\subsection{环节2A：人物卡猜口令（人类直觉版）}
\paragraph{目的}
让学生体验``个性化猜测''：当攻击者知道一些个人信息时，猜口令会比纯暴力更``聪明''。

\paragraph{原理（直观版）}
很多人会把口令写成``有意义的东西''：姓名/昵称、生日、学校缩写、喜欢的球星/游戏、常用数字等。攻击者会优先尝试这些``最可能''的组合。

\paragraph{步骤}
\begin{enumerate}
  \item 每组领取2--3张人物卡（虚构信息），每张卡对应一个``目标账号''（仅课堂使用）。
  \item 对每个角色，写出Top20候选口令并排序（最可能$\rightarrow$最不可能），同时标注模式（如：姓名拼音+生日）。
  \item 全班汇总：把共通模式写在白板上，讨论哪些模式最常见、为什么。
\end{enumerate}

\paragraph{小组输出}
Top20清单（含排序理由）+ 共通模式清单（至少5条）。

\subsection{环节2B：TarGuess-i（PCFG-based）运行与对比}
\paragraph{目的}
让学生看到``统计模型''如何把``模式''变成\textbf{概率排序}，并在同样信息下生成更系统的候选列表。

\paragraph{原理（课堂版，不讲公式）}
\begin{itemize}
  \item TarGuess-i把口令拆成``结构''（例如：字母段+数字段+符号）与``填空''（把姓名、生日、爱好等填入结构）。
  \item 它从大量真实口令中学习：哪些结构更常见、哪些组合更常被人使用，于是得到一个``更可能先试''的顺序。
  \item 所以它不是随机试，而是\textbf{先试最可能的那批}。
\end{itemize}

\paragraph{实验准备（教师）}
\begin{itemize}
  \item 已配置好的TarGuess-i代码与运行脚本（你已准备）。
  \item 人物卡信息对应的输入格式（例如JSON/CSV/文本模板）。
  \item 预先为每张人物卡设置一个``目标口令''（仅用于课堂测评），并确保口令遵循可解释的习惯（而非纯随机）。
\end{itemize}

\paragraph{步骤（学生跟做，教师带跑）}
\begin{enumerate}
  \item 学生把人物卡信息录入TarGuess-i所需格式（或直接使用教师提供的输入文件）。
  \item 运行TarGuess-i生成Top-$N$候选（建议$N=100$或$500$，按课堂节奏调整）。
  \item 查找目标口令在候选中的排名（第几位），并记录。
  \item 对比分析：TarGuess-i的Top20与人类Top20有多少重合？它新增了哪些模式？
\end{enumerate}

\paragraph{对比分析表（Day 2）}
\begin{longtable}{@{}p{0.18\textwidth}p{0.34\textwidth}p{0.38\textwidth}@{}}
\toprule
\textbf{项目} & \textbf{人类Top20} & \textbf{TarGuess-i Top-$N$} \\
\midrule
\endhead
目标口令是否出现 & 是 / 否 & 是 / 否 \\
若出现，排名 & \underline{\hspace{10em}} & \underline{\hspace{10em}} \\
主要模式（举3条） & \underline{\hspace{16em}} & \underline{\hspace{16em}} \\
新增/意外模式（举2条） & --- & \underline{\hspace{16em}} \\
结论一句话 & \multicolumn{2}{p{0.72\textwidth}}{\underline{\hspace{30em}}} \\
\bottomrule
\end{longtable}

\subsection{环节2C：简提LLM与风险升级（10--20分钟）}
\paragraph{讲解要点（只讲趋势）}
\begin{itemize}
  \item 传统统计模型擅长``结构+概率''；大模型可能更擅长从公开文本中学会``更像人写的表达''与梗。
  \item 结论回到防御：不要把口令和可公开的个人信息强绑定；重要系统启用MFA；系统侧要做限速、锁定、风控与监控。
\end{itemize}

\paragraph{课堂互动（可选，10分钟）}
\begin{itemize}
  \item 问题1：如果一个人的公开信息越多，个性化猜测会变得更容易还是更难？为什么？
  \item 问题2：``为了好记，把口令设成自己昵称+生日''有什么问题？给出更好的替代方案。
  \item 问题3：当系统强制很复杂的口令规则时，会带来哪些可用性问题？如何折中（例如密码管理器、MFA）？
\end{itemize}

\subsection{Day 2 汇报与展示（建议20--30分钟）}
\begin{itemize}
  \item 每组用1页纸或1张幻灯片汇报：Top20模式总结、TarGuess-i命中排名、``我会如何改进口令策略/登录流程''。
  \item 教师点评维度：模式是否合理、解释是否清晰、建议是否可执行、是否体现安全与可用性的权衡。
\end{itemize}

\subsection*{Day 2 过程考核表（教师/助教）}
\begin{longtable}{@{}p{0.28\textwidth}p{0.10\textwidth}p{0.54\textwidth}@{}}
\toprule
\textbf{考核点} & \textbf{达成} & \textbf{备注} \\
\midrule
\endhead
能总结至少5条常见口令模式并举例 & $\square$ & \\
完成人物卡Top20排序且理由可解释 & $\square$ & \\
能正确运行TarGuess-i并找到目标口令排名 & $\square$ & \\
能解释TarGuess-i与纯暴力的差别（排序/概率） & $\square$ & \\
能提出2条个人建议+2条系统建议（LLM时代） & $\square$ & \\
\bottomrule
\end{longtable}

\section{评分建议（可直接落到过程考核分）}
\subsection*{小组得分（建议60分）}
\begin{longtable}{@{}p{0.46\textwidth}p{0.12\textwidth}p{0.34\textwidth}@{}}
\toprule
\textbf{项目} & \textbf{分值} & \textbf{评分要点} \\
\midrule
\endhead
实验1观察记录完整 & 10 & 字段齐全、能解释空间与速率 \\
防护改造实现与演示 & 25 & 至少1项可运行；触发条件清晰；提示信息合理 \\
复测对比与结论 & 15 & 前后对比数据可信；能说清楚为何变慢/变难 \\
Day2对比分析（人类vs TarGuess） & 10 & 有Top20、排名、模式总结与反思 \\
\bottomrule
\end{longtable}

\subsection*{个人得分（建议40分）}
\begin{longtable}{@{}p{0.46\textwidth}p{0.12\textwidth}p{0.34\textwidth}@{}}
\toprule
\textbf{项目} & \textbf{分值} & \textbf{评分要点} \\
\midrule
\endhead
课堂参与与分工履职 & 10 & 有实际贡献（开发/测试/记录/汇报） \\
概念理解（口头或小测） & 20 & 能解释HTTPS/哈希/在线爆破/防护逻辑 \\
反思与建议 & 10 & 能给出可执行的个人与系统建议 \\
\bottomrule
\end{longtable}

\section{过程考核表（可打印总表）}
\subsection*{E0：助教巡检记录（每组一张）}
\begin{longtable}{@{}p{0.10\textwidth}p{0.18\textwidth}p{0.22\textwidth}p{0.40\textwidth}@{}}
\toprule
\textbf{时间} & \textbf{巡检对象} & \textbf{当前进度} & \textbf{问题与指导建议} \\
\midrule
\endhead
\underline{\hspace{6em}} & 第\underline{\hspace{2em}}组 & $\square$理解完成 $\square$编码中 $\square$自测中 $\square$待复测 & \underline{\hspace{20em}} \\
\underline{\hspace{6em}} & 第\underline{\hspace{2em}}组 & $\square$理解完成 $\square$编码中 $\square$自测中 $\square$待复测 & \underline{\hspace{20em}} \\
\underline{\hspace{6em}} & 第\underline{\hspace{2em}}组 & $\square$理解完成 $\square$编码中 $\square$自测中 $\square$待复测 & \underline{\hspace{20em}} \\
\bottomrule
\end{longtable}

\subsection*{E1：小组汇报评分表（每组一张）}
\begin{longtable}{@{}p{0.52\textwidth}p{0.14\textwidth}p{0.24\textwidth}@{}}
\toprule
\textbf{维度} & \textbf{分值} & \textbf{得分} \\
\midrule
\endhead
表达清晰：能用通俗语言讲清楚做了什么与为什么有效 & 10 & \underline{\hspace{6em}} \\
证据充分：有前后对比数据/现象支撑结论 & 10 & \underline{\hspace{6em}} \\
工程性：触发条件清晰、提示信息合理、考虑副作用 & 10 & \underline{\hspace{6em}} \\
反思深度：能指出至少1个仍存在的风险与改进方向 & 10 & \underline{\hspace{6em}} \\
\bottomrule
\end{longtable}

\subsection*{E2：个人参与与反思（每人一张，课后5分钟）}
\begin{itemize}
  \item 我在小组中的角色：$\square$开发 $\square$测试 $\square$记录 $\square$汇报 $\square$组长/协调
  \item 我本次贡献的具体工作（写一句话）：\underline{\hspace{24em}}
  \item 我学到的最重要的1个概念：\underline{\hspace{24em}}
  \item 我对自己的账号口令习惯准备做的1个改变：\underline{\hspace{24em}}
  \item 我对网站登录系统的1条改进建议：\underline{\hspace{24em}}
\end{itemize}

\section{（扩展闲谈）离线攻击：什么时候会发生？}
\paragraph{一句话解释}
\textbf{离线攻击}是攻击者拿到了口令的``校验物''（通常是哈希）后，在自己电脑上无限次猜测；此时验证码/Cloudflare对他几乎无效，因为他不再访问你的网站。

\paragraph{课堂可讲的判断法}
如果攻击者\textbf{只有登录页面}、只能对\texttt{/login}发请求：这是\textbf{在线}。\\
如果攻击者拿到了\textbf{数据库备份/用户表哈希/日志泄露的哈希}：才可能做\textbf{离线}。

\section{教师备课清单（建议）}
\begin{itemize}
  \item 练习站点可一键启动；所有账号与口令均为虚构，且口令空间受控（课堂安全）。
  \item 在线尝试演示脚本可控速率、可控次数，带``停止开关''。
  \item 预设人物卡$\times$若干套（建议至少6种画像，便于轮换防止泄题）。
  \item TarGuess-i运行环境已验证：输入$\rightarrow$输出稳定；准备一份``标准输出样例''以防现场故障。
  \item 打印：实验1观察表、实验2对比表、Day2对比分析表、过程考核表。
\end{itemize}

\section{课堂安全与合规声明（建议投影并由学生确认）}
\begin{itemize}
  \item 我承诺：只在老师提供的课堂靶场上进行测试，不测试外网网站与他人账号。
  \item 我承诺：不传播、不保存真实账号口令；所有实验账号均为虚构或一次性。
  \item 我理解：本课程目的是学习防护设计与安全意识，任何超出课堂授权的测试都可能违法违规。
\end{itemize}

\appendix
\section{附录A：人物卡模板（示例，可自行扩展）}
\subsection*{人物卡1（示例）}
\begin{itemize}
  \item 昵称：LinLin\quad 真实姓名（可选）：林某某
  \item 年级：初二\quad 学校简称：THS
  \item 生日：2011-08-16\quad 喜欢的数字：16、8
  \item 爱好：篮球\quad 喜欢的球星：Kobe
  \item 常用ID：linlin16\quad 常用口头禅：``冲冲冲''
  \item （教师私密）目标口令：\underline{\hspace{10em}}
\end{itemize}

\subsection*{人物卡2（示例）}
\begin{itemize}
  \item 昵称：Ming\quad 学校简称：BJ
  \item 生日：2010-12-01\quad 爱好：原神\quad 常用ID：ming\_paimon
  \item 喜欢的数字：1201\quad 喜欢的词：``wind''
  \item （教师私密）目标口令：\underline{\hspace{10em}}
\end{itemize}

\subsection*{人物卡3（示例）}
\begin{itemize}
  \item 昵称：Zoe\quad 学校简称：QHFZ
  \item 生日：2009-03-07\quad 爱好：摄影\quad 口头禅：``今天也要出片''
  \item 喜欢的数字：37\quad 常用ID：zoe\_photo37
  \item 喜欢的词：``light''、``lens''
  \item （教师私密）目标口令：\underline{\hspace{10em}}
\end{itemize}

\subsection*{人物卡4（示例）}
\begin{itemize}
  \item 昵称：Kai\quad 学校简称：BJ
  \item 生日：2011-11-11\quad 爱好：街舞\quad 喜欢的组合：``KDA''
  \item 喜欢的数字：1111\quad 常用ID：kai\_dance11
  \item （教师私密）目标口令：\underline{\hspace{10em}}
\end{itemize}

\subsection*{人物卡5（示例）}
\begin{itemize}
  \item 昵称：HaoHao\quad 学校简称：THS
  \item 生日：2010-06-18\quad 爱好：编程\quad 常用ID：haohack618（注意：虚构）
  \item 喜欢的词：``debug''\quad 喜欢的数字：618
  \item （教师私密）目标口令：\underline{\hspace{10em}}
\end{itemize}

\subsection*{人物卡6（示例）}
\begin{itemize}
  \item 昵称：Yuki\quad 学校简称：QAX
  \item 生日：2008-09-09\quad 爱好：二次元\quad 喜欢的角色名：``Asuna''
  \item 常用ID：yuki\_asuna\_99\quad 喜欢的数字：99
  \item （教师私密）目标口令：\underline{\hspace{10em}}
\end{itemize}

\section{附录B：口令模式清单（课堂引导用）}
可引导学生从下列方向总结模式（避免直接给答案，可先让学生总结后再补充）：
\begin{itemize}
  \item 姓名/昵称拼音 + 年份/生日（全/后两位）
  \item 学校缩写 + 数字（年级/班级/生日片段）
  \item 爱好/偶像 + 数字（球衣号/纪念年份）
  \item ID复用（把用户名直接当口令，或轻微变体）
  \item ``词+词+数字''、``词+符号+数字''（增加可记忆性）
\end{itemize}

\section{附录C：可压缩版时间安排（示例）}
若每一天仅有3小时，可按以下压缩：
\begin{itemize}
  \item Day1：科普30min + 在线爆破演示20min + 防护改造90min + 复测与总结40min
  \item Day2：人物卡猜测40min + TarGuess-i对比70min + LLM/离线闲谈20min + 汇报与评价50min
\end{itemize}

\section{附录D：工作纸模板（可直接打印）}
\subsection*{D1：人物卡Top20候选表}
\begin{longtable}{@{}p{0.06\textwidth}p{0.38\textwidth}p{0.46\textwidth}@{}}
\toprule
\textbf{\#} & \textbf{候选口令（示例写法）} & \textbf{理由/模式标注（如：昵称+生日后四位）} \\
\midrule
\endhead
1 & \underline{\hspace{16em}} & \underline{\hspace{18em}} \\
2 & \underline{\hspace{16em}} & \underline{\hspace{18em}} \\
3 & \underline{\hspace{16em}} & \underline{\hspace{18em}} \\
4 & \underline{\hspace{16em}} & \underline{\hspace{18em}} \\
5 & \underline{\hspace{16em}} & \underline{\hspace{18em}} \\
6 & \underline{\hspace{16em}} & \underline{\hspace{18em}} \\
7 & \underline{\hspace{16em}} & \underline{\hspace{18em}} \\
8 & \underline{\hspace{16em}} & \underline{\hspace{18em}} \\
9 & \underline{\hspace{16em}} & \underline{\hspace{18em}} \\
10 & \underline{\hspace{16em}} & \underline{\hspace{18em}} \\
11 & \underline{\hspace{16em}} & \underline{\hspace{18em}} \\
12 & \underline{\hspace{16em}} & \underline{\hspace{18em}} \\
13 & \underline{\hspace{16em}} & \underline{\hspace{18em}} \\
14 & \underline{\hspace{16em}} & \underline{\hspace{18em}} \\
15 & \underline{\hspace{16em}} & \underline{\hspace{18em}} \\
16 & \underline{\hspace{16em}} & \underline{\hspace{18em}} \\
17 & \underline{\hspace{16em}} & \underline{\hspace{18em}} \\
18 & \underline{\hspace{16em}} & \underline{\hspace{18em}} \\
19 & \underline{\hspace{16em}} & \underline{\hspace{18em}} \\
20 & \underline{\hspace{16em}} & \underline{\hspace{18em}} \\
\bottomrule
\end{longtable}

\subsection*{D2：防护设计说明（1页模板）}
\begin{itemize}
  \item 我们选择的防护（勾选）：$\square$锁定 $\square$延迟/退避 $\square$验证码 $\square$简化2FA $\square$其他：\underline{\hspace{10em}}
  \item 要对抗的攻击特征（至少1条）：\underline{\hspace{24em}}
  \item 触发条件（阈值$N$、时长$T$等）：\underline{\hspace{24em}}
  \item 对正常用户的影响（可用性副作用）：\underline{\hspace{24em}}
  \item 我们的复测结果一句话：\underline{\hspace{24em}}
\end{itemize}

\end{document}
